\documentclass[11pt]{article}
\usepackage[margin=1in]{geometry}
\usepackage{listings}
\usepackage{xcolor}
\usepackage{enumitem}
\usepackage{amsmath}

\lstset{
  language=Java,
  basicstyle=\ttfamily\small,
  keywordstyle=\bfseries\color{blue!70!black},
  commentstyle=\color{green!50!black},
  stringstyle=\color{red!60!black},
  showstringspaces=false,
  frame=single,
  framesep=6pt,
  xleftmargin=6pt,
  xrightmargin=6pt,
  aboveskip=10pt,
  belowskip=10pt,
}

\pagestyle{empty}

\begin{document}

\begin{center}
  {\Large \textbf{CSCI-UA 102 --- Quiz 3}} \\[6pt]
  {\large Recursion} \\[4pt]
  Name: \underline{\hspace{2.5in}} \quad NetID: \underline{\hspace{1.5in}}
\end{center}

\bigskip

\section*{funNum}

For an integer array of size $N$, I define the \textbf{funNum} of the array as follows:

\begin{itemize}[nosep]
  \item If the array has \textbf{no elements}, then funNum is \textbf{0}.
  \item If the array has \textbf{one element}, then funNum is equal to \textbf{that element}.
  \item Otherwise, funNum is equal to the \textbf{sum} of:
    \begin{enumerate}[nosep]
      \item the last element of the array,
      \item the funNum of the first $N-1$ elements, and
      \item the funNum of the first $N-2$ elements.
    \end{enumerate}
\end{itemize}

\bigskip

\noindent\textbf{Examples:}
\begin{itemize}[nosep]
  \item \texttt{funNum(\{\})} $= 0$
  \item \texttt{funNum(\{5\})} $= 5$
  \item \texttt{funNum(\{5, 3\})} $= 3 + 5 + 0 = 8$
  \item \texttt{funNum(\{5, 3, 1\})} $= 1 + 8 + 5 = 14$
\end{itemize}

\bigskip

\noindent\textbf{Write an efficient, recursive method:}
\begin{lstlisting}
public int funNum(int[] array)
\end{lstlisting}

\vfill

\noindent\textbf{What is the time complexity of your solution?} \underline{\hspace{2in}}

\newpage

{\color{blue}
\section*{Solution}

\textbf{Key insight:} \texttt{funNum} only depends on the prefix length, so the recurrence is
$f(n) = \texttt{array}[n{-}1] + f(n{-}1) + f(n{-}2)$, identical in structure to Fibonacci.
Naive recursion is $O(2^n)$; an efficient solution avoids recomputation.

\textbf{Efficient solution --- return previous values} (same pattern as Fibonacci):

\begin{lstlisting}[basicstyle=\ttfamily\small\color{blue},
  keywordstyle=\bfseries\color{blue!70!black},
  commentstyle=\color{blue!50}]
public int funNum(int[] array) {
    if (array.length == 0) return 0;
    if (array.length == 1) return array[0];
    return funNumPair(array, array.length)[0];
}

// returns {funNum(first n elems), funNum(first n-1 elems)}
private int[] funNumPair(int[] array, int n) {
    if (n == 1) return new int[]{array[0], 0};
    int[] prevPair = funNumPair(array, n - 1);
    // prevPair[0] = f(n-1),  prevPair[1] = f(n-2)
    int funNumN = array[n - 1] + prevPair[0] + prevPair[1];
    return new int[]{funNumN, prevPair[0]};
}
\end{lstlisting}

The \texttt{funNumPair} method recurses naturally on $n{-}1$ but returns \emph{two} values: 
$f(n)$ and $f(n{-}1),$ so the caller gets both values it needs from a
single recursive call instead of two, eliminating the branching.

\bigskip
\textbf{Time complexity:} $O(n)$ --- each prefix length is computed exactly once.

\end{document}
